\SourceLeaf{406}%
% A544-P0406-D01
\pdfbookmark[1]{V.}{complexe-V}
\begin{center}[V.]\end{center}
Sollte sich bestätigen, dass jede Einheitszahl blos aus Factoren von der Form
\[
\frac{\varepsilon^\alpha-\varepsilon^\beta}{\varepsilon^\gamma-\varepsilon^\delta}
\]
zusammengesetzt wäre, so würde folgender Satz bewiesen sein:

\emph{Ist $f(\varepsilon)$ eine Einheitszahl, so ist}
\[
\frac{f(\varepsilon)}{f(\varepsilon^{-1})}=\varepsilon^n
\]

% A544-P0406-D02
Auch ohne \emph{jenen} Satz \emph{vorauszusetzen} ist der Schlusssatz leicht zu beweisen.
Es sei
\[
\frac{f\varepsilon}{f\varepsilon^{-1}}=F\varepsilon
\]
so ist
\[
F\varepsilon\pdot F\varepsilon^{-1}=1
\]
woraus mit Hülfe der Lehre von der Mensur leicht gefolgert wird, dass
\[
F\varepsilon=\pm\varepsilon^n
\]
Das untere Zeichen ist aber unmöglich, weil sonst $f\varepsilon$ durch $1-\varepsilon$ theilbar sein müsste.

% A544-P0406-D03
Dass der Determinant einer von $0$ verschiedenen Zahl nicht $=0$ sein könne, lässt sich leicht beweisen. Wenn der Determinant durch $m$ theilbar ist, so ist die Zahl selbst durch $1-\varepsilon$ theilbar; folglich wenn der Determinant durch $m^{m-1}$ theilbar ist, muss die Zahl selbst durch $m$ theilbar sein. Welches absurd ist, da beim Det. $0$ die Zahl erst salvo Det. so oft durch $m$ dividirt werden könnte, bis sie nicht mehr theilbar wäre. Der erste Satz aber erhellt so. Es sei die vorgegebene Zahl
\[
a+b\varepsilon+c\varepsilon\varepsilon+\text{etc.}\equiv a+b+c\dotstwo\text{ mod. }1-\varepsilon
\]
also Determinans $\equiv(a+b+c\dotstwo)^{m-1}\text{ mod. }1-\varepsilon$.
\BookRule{.3\textwidth}

\SourceLeaf{407}%
% A544-P0407-D01
\pdfbookmark[1]{VI.}{complexe-VI}
\begin{center}[VI.]\end{center}
Es sei $\varepsilon^n=1$
\[
\begin{aligned}
f\varepsilon&=a+b\varepsilon+c\varepsilon\varepsilon+d\varepsilon^3+\text{etc.}\\
m&=\text{Determinans dieser Zahl}\\
\frac{m}{f\varepsilon}&=f\varepsilon\varepsilon\pdot f\varepsilon^3\ldots f\varepsilon^{n-1}\\
&=A+B\varepsilon+C\varepsilon\varepsilon+\text{etc.}=F\varepsilon
\end{aligned}
\]
Der Zahl $f\varepsilon$ entspricht eine Wurzel der Congruenz $x^n\equiv1\ (\text{mod. }m.)$ Es sei dieselbe $r$. Man hat
\[
\begin{aligned}
nA&=F1+F\varepsilon+F\varepsilon\varepsilon+\dotstwo\\
nB&=F1+\varepsilon^{-1}F\varepsilon+\varepsilon^{-2}F\varepsilon\varepsilon+\dotstwo\\
nC&=F1+\varepsilon^{-2}F\varepsilon+\varepsilon^{-4}F\varepsilon\varepsilon+\dotstwo\\
&\text{etc.}
\end{aligned}
\]

% A544-P0407-D02
also, da $F\varepsilon\varepsilon,F\varepsilon^3,F\varepsilon^4$ etc. durch $f\varepsilon$ theilbar sind,
\[
\begin{gathered}
nA-F1-\varepsilon(nB-F1)\\
nA-F1-\varepsilon\varepsilon(nC-F1)\\
nA-F1-\varepsilon^3(nD-F1)\\
\text{etc.}
\end{gathered}
\]
alle durch $f\varepsilon$ theilbar, oder auch
\[
\begin{gathered}
n(A-B)-\varepsilon n(B-C)\\
n(B-C)-\varepsilon n(C-D)\\
\text{etc.}
\end{gathered}
\]
durch $f\varepsilon$ theilbar; folglich [wenn $f\varepsilon$ durch $1-\varepsilon$, und $F\varepsilon$ durch eine ganze reelle Zahl nicht theilbar ist]
\[
\varepsilon\equiv\frac{A-B}{B-C}\equiv\frac{B-C}{C-D}\equiv\frac{C-D}{D-E}\text{ etc. }(\text{mod. }f\varepsilon)
\]
\BookRule{.3\textwidth}

\SourceLeaf{408}%
% A544-P0408-D01
\pdfbookmark[1]{Bemerkungen zu den complexen Zahlen}{complexe-Bemerkungen}
\begin{center}BEMERKUNGEN.\end{center}
\BookRule{.15\textwidth}
Die hier unter der gemeinsamen Ueberschrift, zur Theorie der complexen Zahlen, zusammengestellten Untersuchungen bilden zerstreute Notizen in der Handschrift. Sie enthalten die wesentlichen Momente des Beweises vom \textsc{Fermat}schen Satze für die dritte und fünfte Potenz. Die aus dritten Wurzeln der Einheit zusammengesetzten Zahlen sind in unvollständigen hier nicht abgedruckten Aufzeichnungen sowohl mit Hülfe der Theorie der binären quadratischen Formen, als auch der Kreistheilung untersucht. Bei Gelegenheit der Anwendung der letztern und zwar während der Ausarbeitung der Abhandlung \emph{Disquisitionum circa aequationes puras ulterior evolutio} ist noch die ternäre cubische Form aufgestellt, in welche $27\frac{x^n-1}{x-1}$ für eine Primzahl $n\equiv1\text{ mod. }3$ verwandelt werden kann, und zugleich die Theorie der Composition der mit jener verwandten Form $X^3+mY^3+mmZ^3-3mXYZ$ entwickelt.

% A544-P0408-D02
Die in den Untersuchungen des Bruchstück [I] vorausgesetzte Eigenschaft der aus dritten Wurzeln der Einheit gebildeten ganzen Zahlen, dass jede nur auf Eine Weise in Primfactoren zerlegt werden kann, ergibt sich aus dem \textsc{Euclid}ischen Verfahren, die gemeinsamen Theiler zweier Zahlen zu bestimmen, wenn dabei der unter [II] abgeleitete Satz über die nächste ganze Zahl für irgend eine vorgegebene Bruchzahl in Anwendung gebracht wird.

% A544-P0408-D03
Dass dieselbe Fundamentaleigenschaft auch den aus fünften Wurzeln der Einheit zusammengesetzten Zahlen zukommt, folgt daraus, dass der nach einer ganz analogen Regel wie in [II] gebildete Bruchrest entweder von $m$ oder doch von $m$ multiplicirt in eine geeignete Einheitszahl $E$ so beschaffen ist, dass er durch Subtraction von der vorgegebenen Zahl $mE$ eine ganze Zahl entstehen lässt und dass sein Determinant die Einheit nicht übertrifft. Die Einheitszahlen lassen sich aber, wie in [III] angedeutet, aus der Theorie der binären quadratischen Formen vom Determinant $5$ in Verbindung mit der Zerlegung irgend einer reellen Primzahl in vier Factoren (z. B. $11=\text{Det. }(2+\varepsilon)$) ableiten, nemlich als Producte der Potenzen von $\varepsilon$ und $1+\varepsilon$.
\begin{flushright}\textsc{Schering.}\end{flushright}
\BookRule{.3\textwidth}

\clearpage
\SourceLeaf{409}%
% A544-P0409-D01
\renewcommand{\BookRunningTitle}{Tafel des quadratischen Characters}
\pdfbookmark[0]{Tafel des quadratischen Characters}{quadratic-character-table}
\begin{center}
{\Large TAFEL}\par\bigskip
DES QUADRATISCHEN CHARACTERS\par\bigskip
DER PRIMZAHLEN VON 2 BIS 997 ALS RESTE\par\bigskip
IN BEZUG\par\bigskip
AUF DIE PRIMZAHLEN VON 3 BIS 503 ALS THEILER.
\end{center}
\clearpage

\SourceLeaf{410}%
% A544-P0410-D01
\begin{center}\begingroup\fontsize{8}{10.5}\selectfont\setlength{\tabcolsep}{2pt}
\begin{tabular}{r|c|c|c|c|c|c|c|c|c|c|c|c|c|c|c|c|c|r}
\hline
 & 2 & 3 & 5 & 7 & 11 & 13 & 17 & 19 & 23 & 29 & 31 & 37 & 41 & 43 & 47 & 53 & 59 & \\
\hline
3 &  & $-$ &  & $-$ &  & $-$ &  & $-$ &  &  & $-$ & $-$ &  & $-$ &  &  &  & 3\\
5 &  &  & $-$ &  & $-$ &  &  & $-$ &  & $-$ & $-$ &  & $-$ &  &  &  & $-$ & 5\\
7 & $-$ &  &  & $-$ & $-$ &  &  &  & $-$ & $-$ &  & $-$ &  & $-$ &  & $-$ &  & 7\\
11 &  & $-$ & $-$ &  & $-$ &  &  &  & $-$ &  & $-$ & $-$ &  &  & $-$ & $-$ & $-$ & 11\\
13 &  & $-$ &  &  &  & $-$ & $-$ &  & $-$ & $-$ &  &  &  & $-$ &  & $-$ &  & 13\\
17 & $-$ &  &  &  &  & $-$ & $-$ & $-$ &  &  &  &  &  & $-$ & $-$ & $-$ & $-$ & 17\\
19 &  &  & $-$ & $-$ & $-$ &  & $-$ & $-$ & $-$ &  &  &  &  & $-$ & $-$ &  &  & 19\\
23 & $-$ & $-$ &  &  &  & $-$ &  &  & $-$ & $-$ & $-$ &  & $-$ &  & $-$ &  & $-$ & 23\\
29 &  &  & $-$ & $-$ &  & $-$ &  &  & $-$ & $-$ &  &  &  &  &  & $-$ & $-$ & 29\\
31 & $-$ &  & $-$ & $-$ &  &  &  & $-$ &  &  & $-$ &  & $-$ &  & $-$ &  & $-$ & 31\\
37 &  & $-$ &  & $-$ & $-$ &  &  &  &  &  &  & $-$ & $-$ &  & $-$ & $-$ &  & 37\\
41 & $-$ &  & $-$ &  &  &  &  &  & $-$ &  & $-$ & $-$ & $-$ & $-$ &  &  & $-$ & 41\\
43 &  &  &  &  & $-$ & $-$ & $-$ &  & $-$ &  & $-$ &  & $-$ & $-$ & $-$ & $-$ & $-$ & 43\\
47 & $-$ & $-$ &  & $-$ &  &  & $-$ &  &  &  &  & $-$ &  &  & $-$ & $-$ & $-$ & 47\\
53 &  &  &  & $-$ & $-$ & $-$ & $-$ &  &  & $-$ &  & $-$ &  & $-$ & $-$ & $-$ & $-$ & 53\\
59 &  & $-$ & $-$ & $-$ &  &  & $-$ & $-$ &  & $-$ &  &  & $-$ &  &  & $-$ & $-$ & 59\\
61 &  & $-$ & $-$ &  &  & $-$ &  & $-$ &  &  &  &  & $-$ &  & $-$ &  &  & 61\\
67 &  &  &  &  &  &  & $-$ & $-$ & $-$ & $-$ &  & $-$ &  &  & $-$ &  & $-$ & 67\\
71 & $-$ & $-$ & $-$ &  &  &  &  & $-$ &  & $-$ &  & $-$ &  & $-$ &  &  &  & 71\\
73 & $-$ & $-$ &  &  &  &  &  & $-$ & $-$ &  &  & $-$ & $-$ &  &  &  &  & 73\\
79 & $-$ &  & $-$ &  & $-$ & $-$ &  & $-$ & $-$ &  & $-$ &  &  &  &  &  &  & 79\\
83 &  & $-$ &  & $-$ & $-$ &  & $-$ &  & $-$ & $-$ & $-$ & $-$ & $-$ &  &  &  & $-$ & 83\\
89 & $-$ &  & $-$ &  & $-$ &  & $-$ &  &  &  &  &  &  &  & $-$ & $-$ &  & 89\\
97 & $-$ & $-$ &  &  & $-$ &  &  &  &  &  & $-$ &  &  & $-$ & $-$ & $-$ &  & 97\\
101 &  &  & $-$ &  &  & $-$ & $-$ & $-$ & $-$ &  & $-$ & $-$ &  & $-$ & $-$ &  &  & 101\\
103 & $-$ &  &  & $-$ &  & $-$ & $-$ & $-$ & $-$ & $-$ &  &  & $-$ &  &  &  & $-$ & 103\\
107 &  & $-$ &  &  & $-$ & $-$ &  & $-$ & $-$ & $-$ &  & $-$ & $-$ &  & $-$ & $-$ &  & 107\\
109 &  & $-$ & $-$ & $-$ &  &  &  &  &  & $-$ & $-$ &  &  & $-$ &  &  &  & 109\\
113 & $-$ &  &  & $-$ & $-$ & $-$ &  &  &  &  & $-$ &  & $-$ &  &  & $-$ &  & 113\\
127 & $-$ &  &  &  & $-$ & $-$ & $-$ & $-$ &  &  & $-$ & $-$ & $-$ &  & $-$ &  &  & 127\\
131 &  & $-$ & $-$ & $-$ & $-$ & $-$ &  &  &  &  &  &  & $-$ & $-$ &  & $-$ & $-$ & 131\\
137 & $-$ &  &  & $-$ & $-$ &  & $-$ & $-$ &  &  &  & $-$ &  &  &  &  & $-$ & 137\\
139 &  &  & $-$ & $-$ & $-$ & $-$ &  &  &  & $-$ & $-$ & $-$ & $-$ &  & $-$ &  &  & 139\\
149 &  &  & $-$ & $-$ &  &  & $-$ & $-$ &  & $-$ & $-$ & $-$ &  &  & $-$ & $-$ &  & 149\\
151 & $-$ &  & $-$ &  & $-$ &  & $-$ & $-$ &  & $-$ & $-$ & $-$ &  & $-$ & $-$ &  & $-$ & 151\\
157 &  & $-$ &  &  & $-$ & $-$ & $-$ & $-$ &  &  & $-$ & $-$ &  &  & $-$ &  &  & 157\\
163 &  &  &  &  &  &  &  &  &  &  &  &  & $-$ & $-$ & $-$ & $-$ &  & 163\\
167 & $-$ & $-$ &  & $-$ & $-$ &  &  & $-$ &  & $-$ & $-$ &  &  &  & $-$ &  &  & 167\\
173 &  &  &  &  &  & $-$ &  &  & $-$ & $-$ & $-$ & $-$ & $-$ & $-$ & $-$ &  &  & 173\\
179 &  & $-$ & $-$ &  &  & $-$ & $-$ & $-$ &  & $-$ & $-$ &  &  & $-$ & $-$ &  & $-$ & 179\\
181 &  & $-$ & $-$ &  & $-$ & $-$ &  &  &  & $-$ &  & $-$ &  & $-$ &  &  & $-$ & 181\\
191 & $-$ & $-$ & $-$ &  &  & $-$ & $-$ &  & $-$ &  &  &  &  & $-$ &  &  & $-$ & 191\\
193 & $-$ & $-$ &  & $-$ &  &  &  &  & $-$ &  & $-$ &  &  & $-$ &  &  & $-$ & 193\\
197 &  &  &  & $-$ &  &  &  & $-$ & $-$ & $-$ &  & $-$ & $-$ & $-$ & $-$ & $-$ & $-$ & 197\\
199 & $-$ &  & $-$ & $-$ &  & $-$ &  &  & $-$ & $-$ & $-$ &  &  & $-$ & $-$ & $-$ &  & 199\\
211 &  &  & $-$ &  & $-$ & $-$ &  & $-$ &  &  &  & $-$ &  & $-$ & $-$ & $-$ & $-$ & 211\\
223 & $-$ &  &  & $-$ &  &  & $-$ & $-$ &  & $-$ & $-$ & $-$ & $-$ & $-$ & $-$ & $-$ &  & 223\\
227 &  & $-$ &  & $-$ & $-$ &  &  & $-$ & $-$ & $-$ &  &  &  & $-$ & $-$ & $-$ & $-$ & 227\\
\hline
 & 2 & 3 & 5 & 7 & 11 & 13 & 17 & 19 & 23 & 29 & 31 & 37 & 41 & 43 & 47 & 53 & 59 & \\
\hline
\end{tabular}\endgroup\end{center}
\clearpage
\begin{center}\begingroup\fontsize{8}{10.5}\selectfont\setlength{\tabcolsep}{2pt}
\begin{tabular}{r|c|c|c|c|c|c|c|c|c|c|c|c|c|c|c|c|c|r}
\hline
 & 61 & 67 & 71 & 73 & 79 & 83 & 89 & 97 & 101 & 103 & 107 & 109 & 113 & 127 & 131 & 137 & 139 & \\
\hline
3 & $-$ & $-$ &  & $-$ & $-$ &  &  & $-$ &  & $-$ &  & $-$ &  & $-$ &  &  & $-$ & 3\\
5 & $-$ &  & $-$ &  & $-$ &  & $-$ &  & $-$ &  &  & $-$ &  &  & $-$ &  & $-$ & 5\\
7 &  & $-$ & $-$ &  & $-$ &  &  &  &  &  & $-$ & $-$ & $-$ & $-$ &  & $-$ &  & 7\\
11 &  & $-$ & $-$ &  &  &  & $-$ & $-$ &  & $-$ &  &  & $-$ &  &  & $-$ &  & 11\\
13 & $-$ &  &  &  & $-$ &  &  &  & $-$ & $-$ & $-$ &  & $-$ & $-$ & $-$ &  & $-$ & 13\\
17 &  & $-$ &  &  &  & $-$ & $-$ &  & $-$ & $-$ &  &  &  & $-$ &  & $-$ &  & 17\\
19 & $-$ &  &  & $-$ &  & $-$ &  &  & $-$ &  &  &  &  &  & $-$ & $-$ & $-$ & 19\\
23 &  &  & $-$ & $-$ &  &  &  &  & $-$ &  &  &  &  & $-$ & $-$ &  & $-$ & 23\\
29 &  & $-$ & $-$ &  &  & $-$ &  &  &  & $-$ & $-$ & $-$ &  &  &  &  & $-$ & 29\\
31 &  & $-$ & $-$ &  &  &  &  & $-$ & $-$ & $-$ & $-$ & $-$ & $-$ &  & $-$ &  &  & 31\\
37 &  & $-$ & $-$ & $-$ &  & $-$ &  &  & $-$ &  & $-$ &  &  & $-$ &  & $-$ & $-$ & 37\\
41 & $-$ &  &  & $-$ &  & $-$ &  &  &  & $-$ & $-$ &  & $-$ & $-$ & $-$ &  & $-$ & 41\\
43 &  & $-$ &  &  & $-$ & $-$ &  & $-$ & $-$ & $-$ & $-$ & $-$ &  & $-$ &  &  & $-$ & 43\\
47 & $-$ &  & $-$ &  & $-$ & $-$ & $-$ & $-$ & $-$ & $-$ &  &  &  &  & $-$ &  &  & 47\\
53 &  &  &  &  &  &  & $-$ & $-$ &  &  & $-$ &  & $-$ &  & $-$ &  &  & 53\\
59 &  &  & $-$ &  & $-$ &  &  &  &  &  & $-$ &  &  & $-$ &  & $-$ & $-$ & 59\\
61 & $-$ &  &  & $-$ &  & $-$ &  & $-$ &  & $-$ & $-$ & $-$ & $-$ & $-$ & $-$ & $-$ &  & 61\\
67 &  & $-$ & $-$ & $-$ &  & $-$ & $-$ &  &  & $-$ & $-$ &  &  & $-$ & $-$ &  &  & 67\\
71 &  &  & $-$ & $-$ & $-$ & $-$ & $-$ &  & $-$ & $-$ & $-$ & $-$ &  &  & $-$ &  &  & 71\\
73 & $-$ & $-$ & $-$ & $-$ & $-$ &  & $-$ & $-$ &  &  &  & $-$ &  & $-$ &  & $-$ &  & 73\\
79 &  & $-$ &  & $-$ & $-$ & $-$ & $-$ & $-$ & $-$ &  &  &  &  &  & $-$ &  &  & 79\\
83 & $-$ &  &  &  &  & $-$ &  &  &  &  &  & $-$ & $-$ & $-$ & $-$ &  &  & 83\\
89 &  & $-$ & $-$ & $-$ & $-$ &  & $-$ & $-$ &  &  & $-$ & $-$ &  &  & $-$ &  & $-$ & 89\\
97 & $-$ &  &  & $-$ & $-$ &  & $-$ & $-$ & $-$ & $-$ &  & $-$ & $-$ &  &  &  &  & 97\\
101 &  &  & $-$ &  & $-$ &  &  & $-$ & $-$ &  & $-$ &  &  &  & $-$ & $-$ &  & 101\\
103 & $-$ &  &  &  & $-$ & $-$ &  & $-$ &  & $-$ & $-$ &  &  &  & $-$ & $-$ & $-$ & 103\\
107 & $-$ &  &  &  & $-$ & $-$ & $-$ &  & $-$ &  & $-$ &  &  &  &  & $-$ &  & 107\\
109 & $-$ &  & $-$ & $-$ &  & $-$ & $-$ & $-$ &  &  &  & $-$ & $-$ &  & $-$ & $-$ &  & 109\\
113 & $-$ &  &  &  &  & $-$ &  & $-$ &  &  &  & $-$ & $-$ & $-$ & $-$ &  & $-$ & 113\\
127 & $-$ &  & $-$ & $-$ & $-$ &  &  &  &  & $-$ & $-$ &  & $-$ & $-$ & $-$ &  &  & 127\\
131 & $-$ &  &  &  &  &  & $-$ &  & $-$ &  & $-$ & $-$ & $-$ &  & $-$ &  &  & 131\\
137 & $-$ &  &  & $-$ &  &  &  &  & $-$ & $-$ & $-$ & $-$ &  &  &  & $-$ & $-$ & 137\\
139 &  & $-$ & $-$ &  & $-$ & $-$ & $-$ &  &  &  & $-$ &  & $-$ & $-$ & $-$ & $-$ & $-$ & 139\\
149 & $-$ & $-$ &  & $-$ &  &  &  &  &  & $-$ & $-$ &  & $-$ & $-$ &  &  &  & 149\\
151 &  &  &  &  &  &  &  & $-$ &  & $-$ &  &  &  & $-$ &  & $-$ & $-$ & 151\\
157 &  & $-$ & $-$ &  &  &  & $-$ &  & $-$ &  &  & $-$ & $-$ & $-$ &  &  &  & 157\\
163 & $-$ &  & $-$ &  &  & $-$ &  & $-$ &  &  &  &  & $-$ &  & $-$ &  &  & 163\\
167 & $-$ &  &  &  &  &  & $-$ & $-$ &  &  & $-$ &  &  & $-$ &  & $-$ &  & 167\\
173 &  & $-$ &  & $-$ &  & $-$ & $-$ &  &  &  &  & $-$ & $-$ &  &  & $-$ & $-$ & 173\\
179 & $-$ & $-$ &  &  &  & $-$ & $-$ &  & $-$ &  & $-$ &  &  &  &  &  & $-$ & 179\\
181 &  & $-$ &  & $-$ & $-$ &  &  &  & $-$ &  &  &  &  &  &  & $-$ & $-$ & 181\\
191 &  & $-$ &  &  & $-$ &  &  & $-$ &  & $-$ & $-$ & $-$ &  &  &  &  &  & 191\\
193 &  & $-$ &  &  &  & $-$ &  & $-$ & $-$ &  & $-$ & $-$ &  &  & $-$ & $-$ & $-$ & 193\\
197 & $-$ &  &  &  &  & $-$ &  & $-$ & $-$ &  & $-$ & $-$ &  & $-$ &  & $-$ &  & 197\\
199 & $-$ &  &  &  & $-$ &  & $-$ &  &  & $-$ &  &  &  &  & $-$ &  & $-$ & 199\\
211 &  &  & $-$ & $-$ & $-$ & $-$ &  &  & $-$ & $-$ & $-$ & $-$ & $-$ &  &  & $-$ & $-$ & 211\\
223 &  &  &  & $-$ &  & $-$ & $-$ &  & $-$ &  &  & $-$ &  & $-$ & $-$ &  & $-$ & 223\\
227 &  &  & $-$ & $-$ & $-$ &  & $-$ & $-$ & $-$ & $-$ &  & $-$ & $-$ &  & $-$ &  & $-$ & 227\\
\hline
 & 61 & 67 & 71 & 73 & 79 & 83 & 89 & 97 & 101 & 103 & 107 & 109 & 113 & 127 & 131 & 137 & 139 & \\
\hline
\end{tabular}\endgroup\end{center}

